\section{ODE 求解基础}

\subsection{为什么需要快速 ODE 求解器}

扩散模型的采样过程等价于从 $t=T$（纯噪声）到 $t=0$（干净数据）求解 PF-ODE。朴素方法——如 DDPM 的 1000 步采样——本质上是对 SDE 的欧拉-丸山离散化，每步都需要一次神经网络前向传播。实际应用中，我们希望将步数从数百步压缩到 10$\sim$20 步，同时保持图像质量。

\begin{importantbox}{核心挑战：少步高质量采样}
每一次 ODE 求解器的"步"（step）都需要调用一次噪声预测网络 $\boldsymbol{\epsilon}_\theta$，这是计算的瓶颈。因此，减少求解步数 = 减少网络推理次数 = 加速生成。目标是在 10$\sim$20 次网络评估（NFE, Number of Function Evaluations）内生成高质量图像。
\end{importantbox}

\subsection{欧拉法（Euler Method）}

给定 ODE $\frac{d\mathbf{x}}{dt} = \mathbf{v}(\mathbf{x}, t)$，欧拉法是最简单的数值积分方案：
$$\mathbf{x}_{t+h} = \mathbf{x}_t + h \cdot \mathbf{v}(\mathbf{x}_t, t)$$

其中 $h$ 为步长。直觉上，欧拉法假设在 $[t, t+h]$ 区间内速度场 $\mathbf{v}$ 不变，用起点处的速度做线性外推。

\begin{knowledgebox}{欧拉法的几何直觉}
想象一条在二维平面上的曲线（ODE 的真实轨迹）。欧拉法在当前点计算切线方向（速度场），然后沿切线走一小步。步长越小，近似越准确，但需要更多步。当步长较大时，直线近似偏离真实曲线，产生截断误差（truncation error）。对于扩散模型，DDIM 的采样公式实际上就是对 PF-ODE 的一种欧拉离散化。
\end{knowledgebox}

\subsection{高阶 ODE 求解器概览}

为了在较大步长下仍然保持精度，数值分析领域发展了大量高阶求解器：

\begin{table}[H]
\centering
\begin{tabular}{lcc}
\toprule
\textbf{求解器} & \textbf{阶数} & \textbf{每步 NFE} \\
\midrule
Euler（欧拉法） & 1 & 1 \\
Heun（改进欧拉法/梯形法） & 2 & 2 \\
Runge-Kutta 4 (RK4) & 4 & 4 \\
Dormand-Prince (DOPRI5) & 5 & 6 \\
\bottomrule
\end{tabular}
\caption{常见通用 ODE 求解器及其特性}
\end{table}

\begin{warningbox}{通用求解器的局限性}
上述求解器对\textbf{任意} ODE 都适用，但它们没有利用 PF-ODE 的特殊结构。例如 Runge-Kutta 4 每步需要 4 次网络评估，在扩散模型场景下代价昂贵。DPM-Solver 的核心创新在于：利用 PF-ODE 的半线性结构，在\textbf{每步仅 1 次网络评估}的前提下实现高阶精度。
\end{warningbox}

\subsection{本章小结}

欧拉法是最基础的 ODE 求解方案，也是 DDIM 采样的等价形式。通用高阶求解器（如 RK4）可以提高精度，但每步需要多次网络评估。本讲后续的重点是如何利用扩散模型 ODE 的特殊结构来设计"免费午餐"式的高阶求解器。


